\section*{الفصل \ref{ch:b1:lipids} --- الدهون}

\begin{solution}{exo:b1:lipids:1}
\omterm{def:b1:lipids:fattyacid}{الدهن} جزيء حيوي غير منحل في الماء منحل في المذيبات غير القطبية.
والشحوم والفوسفوليبيدات \omterm{def:b1:lipids:amphiphilic}{والستيرولات} والكاروتينويدات لا تشترك في هيكل
واحد، بل في هذا السلوك تجاه الماء وحده، وهو ما يمنحها أدوارها
المشتركة (المخازن، والغلالات، \omterm{def:b1:membranes-transport:membrane}{والأغشية}).
\end{solution}

\begin{solution}{exo:b1:lipids:2}
C18:2 $\Delta^{9,12}$. وآخر رابطة مزدوجة تبدأ عند الكربون 12 من
الكربوكسيل، أي عند الكربون 6 من نهاية الميثيل: فهو أوميغا-6.
\end{solution}

\begin{solution}{exo:b1:lipids:3}
\omterm{def:b1:lipids:triglyceride}{ثلاثي الغليسيريد} لا رأس \omterm{def:b1:water-small-molecules:water}{قطبي} له: فهو كاره للماء كله، ويُستبعد من
الماء طورًا كتليًا، أي قطيرة. أما \omterm{def:b1:lipids:amphiphilic}{الفوسفوليبيد} فمزدوج الألفة: فرأسه
\omterm{def:b1:water-small-molecules:water}{القطبي} يجب أن يبقى في الماء وذيلاه يجب أن يغادراه، وهذا لا تحققه إلا
صفيحة سماكتها جزيئان.
\end{solution}

\begin{solution}{exo:b1:lipids:4}
نحو \qty{41}{\celsius}، و\qty{-5}{\celsius}، و\qty{20}{\celsius}
(فمنحني \omterm{def:b1:lipids:amphiphilic}{الكولستيرول} لا انتقال حادًا له؛ ونصف ارتفاعه قرب
\qty{20}{\celsius}). والبكتيريا عند \qty{10}{\celsius} تحتاج إلى
التركيب \omterm{def:b1:lipids:fattyacid}{غير المشبع}، المائع دون حرارة نموها بكثير.
\end{solution}

\begin{solution}{exo:b1:lipids:5}
$15\,000\times 38 = \qty{570}{MJ}$؛ و$570/8 = 71$ يومًا.
والغليكوجين: جافًا $570\,000/17 = \qty{33.5}{kg}$، ومماهًا
\qty{134}{kg}.
\end{solution}

\begin{solution}{exo:b1:lipids:6}
C22:0 $>$ C18:0 $>$ C14:0 $>$ C18:1 $>$ C18:3. فبين الحموض المشبعة
تعقد السلاسل الأطول تماسات \omterm{def:b1:water-small-molecules:bonds}{فان دير فالس} أكثر؛ ورابطة مزدوجة
\emph{سيس} واحدة تزيل من الارتصاص أكثر مما تضيفه أربع ذرات كربون
إضافية (\omterm{def:b1:water-small-molecules:ph}{فحمض} C18:1 ينصهر دون C14:0)؛ وكل رابطة أخرى تخفضها من جديد.
\end{solution}

\begin{solution}{exo:b1:lipids:7}
$2\times\qty{140}{\micro m^2}/\qty{0.6}{nm^2} = 2\times 1.4\times
10^{-10}/6\times 10^{-19} = \num{4.7e8}$ جزيء؛ والكتلة
$\num{4.7e8}\times 750/\num{6e23} = \qty{5.8e-13}{g} = \qty{0.58}{pg}$.
\end{solution}

\begin{solution}{exo:b1:lipids:8}
الحجم \qty{1.11}{mL} $= \qty{1.11e-6}{m^3}$؛ وسطح الكرات
$= 6V/d = 6\times 1.11\times 10^{-6}/10^{-6} = \qty{6.7}{m^2}$. أما
قطرة واحدة بهذا الحجم فقطرها $d = \qty{1.28}{cm}$ وسطحها
\qty{5.2}{cm^2}: فسطح الفصل في المستحلب أكبر بثلاثة عشر ألف مرة.
واللِّيباز منحل في الماء ولا يعمل إلا عند سطح الفصل؛ والمعدل يتناسب
معه.
\end{solution}

\begin{solution}{exo:b1:lipids:9}
فوق $T_m$ تعوق حلقاته الصلبة حركة السلاسل المجاورة فتخفض \omterm{def:b1:lipids:fluidity}{الميوعة}؛
ودون $T_m$ يمنع حجمه السلاسل من الارتصاص في الهلام المنتظم فيمنع
التجمد. ومع استحالة الارتصاص التعاوني لا توجد حرارة تتغير عندها
الطبقة المضاعفة كلها دفعة واحدة: فينتشر الانتقال على مدى.
\end{solution}

\begin{solution}{exo:b1:lipids:10}
بسلاسل مشبعة كلها: عند \qty{25}{\celsius} تكون \omterm{def:b1:membranes-transport:membrane}{الأغشية} قريبة من
انتقالها أصلًا وقاسية؛ وعند \qty{5}{\celsius} تتهلمن. \omterm{def:b1:membranes-transport:membrane}{والأغشية}
المتهلمنة تجعل البروتينات تكف عن العمل، ثم تتشقق وترشح عند إعادة
التسخين أو عند إجهاد ميكانيكي: فتفقد \omterm{def:b1:cell-unit-of-life:cell}{الخلايا} محتوياتها ويموت \omterm{def:b1:body-plans-tissues:tissue}{النسيج}
(أذية التبريد). أما النمط البري فينزع الإشباع من دهونه في الخريف
فيبقى مائعًا عند \qty{5}{\celsius}.
\end{solution}

\begin{solution}{exo:b1:lipids:11}
\qty{1.07}{kg} من الماء لكل كيلوغرام \omterm{def:b1:lipids:triglyceride}{شحم}. والأكسجين: \qty{2000}{L}؛
وعند استخلاص \qty{5}{\%} من أكسجين نسبته \qty{21}{\%}، يكون الهواء
اللازم $2000/(0.21\times 0.05) = \qty{190}{m^3}$؛ والماء المفقود عند
\qty{30}{g/m^3}: \qty{5.7}{kg}. فالجمل يفقد بالتنفس خمسة أمثال ما
يعطيه \omterm{def:b1:lipids:triglyceride}{الشحم} من ماء: فالسنام مخزن طاقة لا مخزن ماء.
\end{solution}

\begin{solution}{exo:b1:lipids:12}
إذا وُجدت طبقة مضاعفة، جعلها \omterm{prop:b1:water-small-molecules:properties}{الأثر الكاره للماء} تلتئم وتنغلق في
حويصلات وتقبل دهونًا جديدة بلا إدخال طاقة: فالتجمع تلقائي. لكن
\omterm{def:b1:cell-unit-of-life:cell}{الخلية} تصنع دهونها بإنزيمات تقع في غشاء موجود أصلًا (\omterm{def:b1:eukaryotic-cell:endomembrane}{الشبكة
الإندوبلازمية})، وهي تقحمها في مكانها؛ والغشاء الجديد ينمو باتساع
غشاء قديم ويُقتسم عند الانقسام، حتى إن كل غشاء في كل \omterm{def:b1:cell-unit-of-life:cell}{خلية} منحدر من
\omterm{def:b1:membranes-transport:membrane}{أغشية} \omterm{def:b1:cell-unit-of-life:cell}{الخلية} السابقة. \omterm{prop:b1:lipids:selfassembly}{فالتجمع الذاتي} يفسر فيزياء الصفيحة، لا منشأها
في \omterm{def:b1:cell-unit-of-life:cell}{الخلية}.
\end{solution}

\begin{solution}{pb:b1:lipids:1}
\textbf{1.} C16:0 (0)، وC18:1 $\Delta^9$ (1)، وC22:6
$\Delta^{4,7,10,13,16,19}$ (6).
\textbf{2.} في الدفء: $0.40\times 0 + 0.45\times 1 + 0.15\times 6 = 1.35$؛
وفي البرد: $0 + 0.40 + 2.10 = 2.50$ رابطة مزدوجة لكل سلسلة.
\textbf{3.} كل رابطة \emph{سيس} تثني السلسلة وتمنع الارتصاص الوثيق مع
الجارات؛ فتماسات \omterm{def:b1:water-small-molecules:bonds}{فان دير فالس} أقل، والحرارة اللازمة لاضطراب السلاسل
أقل، و$T_m$ أدنى.
\textbf{4.} $1.15$ رابطة إضافية لكل سلسلة: أي أدنى بنحو
\qty{45}{\celsius}.
\textbf{5.} في الهلام لا تستطيع \omterm{def:b1:lipids:fattyacid}{الدهون} أن تتحرك: فلا تستطيع \omterm{def:b1:membranes-transport:transporters}{النواقل}
أن تغير تشكلها، وتتوقف \omterm{def:b1:membranes-transport:transporters}{المضخات}، وتلتصق \omterm{def:b1:membranes-transport:transporters}{القنوات}؛ وتضمحل تدرجات
الشوارد، ويفشل النقل العصبي، وتُشل السمكة وتموت من برد تحتمله سمكة
مؤقلمة على البرد.
\textbf{6.} نازعات الإشباع، التي تُدخل روابط مزدوجة في سلاسل الأسيل
الدسمة؛ وهي بروتينات مندمجة في \omterm{def:b1:eukaryotic-cell:endomembrane}{الشبكة الإندوبلازمية}، حيث تُصنع
\omterm{def:b1:lipids:fattyacid}{الدهون}.
\textbf{7.} في الحالتين يهدّئ \omterm{def:b1:lipids:fluidity}{الميوعة}: فيصلّب الغشاء البارد الشديد
عدم الإشباع فوق $T_m$ المنخفضة له، ويمنع الغشاء الدافئ من التهلمن في
يوم بارد.
\textbf{8.} $30\,000\times 38 = \qty{1140}{MJ}$.
\textbf{9.} $1140/60 = 19$ يومًا.
\textbf{10.} جافًا: $\num{1.14e6}/17 = \qty{67}{kg}$؛ ومماهًا:
\qty{268}{kg}.
\textbf{11.} $268 - 30 = \qty{238}{kg}$، أي نحو نصف كتلة الجمل.
\textbf{12.} طبقة من \omterm{def:b1:lipids:triglyceride}{الشحم} تحت الجلد كله كانت لتعزل الجسم وتمنعه من
طرح الحرارة في الصحراء؛ أما تركيز \omterm{def:b1:lipids:triglyceride}{الشحم} في سنام واحد فيترك بقية
الجلد حرة لتفقد الحرارة.
\textbf{13.} $1.07\times 30 = \qty{32}{kg}$ من الماء.
\textbf{14.} $2.0\times 30\,000 = \qty{60000}{L} = \qty{60}{m^3}$ من
الأكسجين.
\textbf{15.} $60/(0.21\times 0.05) = \qty{5700}{m^3}$ من الهواء.
\textbf{16.} $5700\times 30 = \qty{171}{kg}$ من الماء.
\textbf{17.} المفقود خمسة أمثال المكتسب: فأكسدة السنام تكلف ماءً.
والسنام مخزن طاقة يتيح للجمل أن يمضي بلا غذاء؛ أما اقتصاده في الماء
فمن كليتيه، ومن احتماله التجفاف، ومن قدرته على ترك حرارته ترتفع.
\textbf{18.} $500\times 3.5\times 6 = \qty{10.5}{MJ}$ تُدخر حرارةً وتُطرح ليلًا؛ وتبخيرها كان ليكلف $10\,500/2.4 = \qty{4.4}{kg}$ من الماء في اليوم.
\textbf{19.} لا يستطيع الغشاء أن يتمدد إلا بضعة في المئة، فمضاعفة
حجم كرية حمراء تقتضي فائضًا كبيرًا من الغشاء مطويًا في الشكل المقعر
الوجهين في الراحة، وهيكلًا خلويًا يدعه ينبسط دون أن يتمزق: فكريات
الجمل الحمراء بيضوية، وفيها غشاء لكل حجم أكثر مما في كريات الإنسان.
\textbf{20.} $2\times 80\times 10^{-12}/0.65\times 10^{-18} = \num{2.5e8}$.
\textbf{21.} $\num{2.5e8}\times\num{8e12}\times 40 = \num{7.9e22}$.
\textbf{22.} $\num{7.9e22}\times 750/\num{6e23} = \qty{99}{g}$.
\textbf{23.} الانتشار الجانبي يبقي الرأس في الماء والذيول في \omterm{def:b1:lipids:triglyceride}{الزيت}
عند كل خطوة، فلا يكلف شيئًا؛ أما الانقلاب فيقتضي أن يعبر الرأس \omterm{def:b1:water-small-molecules:water}{القطبي}
اللبَّ الكاره للماء، وهو حاجز طاقي كبير لا يتجاوزه الاضطراب الحراري
إلا مرة في الأسبوع.
\textbf{24.} الفرق بين الصفيحتين لا يضمحل إلا بسرعة انقلاب \omterm{def:b1:lipids:fattyacid}{الدهون}؛
وعند مرة في الأسبوع لكل جزيء، يبقى اللاتناظر الذي تقيمه إنزيمات
تنقل دهونًا بعينها عبر الغشاء (القالبات) قائمًا إلى ما لا نهاية أمام
الانقلاب التلقائي.
\textbf{25.} نحو \qty{240}{kg} موفَّرة --- أي نحو نصف كتلة جسمه ---
بحمل \qty{30}{kg} من \omterm{def:b1:lipids:triglyceride}{الشحم} بدل \qty{270}{kg} من الغليكوجين المماه؛
والسنام ليس مخزن ماء: فأكسدته تفقد في النفَس خمسة أمثال ما تنتجه من
ماء.
\end{solution}
